\documentclass[../../main.tex]{subfiles}% 注意这里的文件路径不能用 ./main.tex ,否则用latexmk编译子文件会报错
\graphicspath{{\subfix{./image/}}{\subfix{../../image/}}} % 指定图片引用路径,后续可以直接使用图片文件名
% 如果图片存放在上级目录,则使用 ../../image/ 作为备用路径

% 例如:
% \begin{figure}[H]
% \centering
% \includegraphics[width=0.3\textwidth]{图.png}
% \caption{}
% \label{figure:图}
% \end{figure}
% 注意:上述\label{}一定要放在\caption{}之后,否则引用图片序号会只会显示??.

\begin{document}

\section{有界变差函数}

\begin{definition}[有界变差函数]
设 $f(x)$ 是定义在 $[a,b]$ 上的实值函数, 作分划 $\Delta$: $a=x_0<x_1<\cdots<x_n=b$ 以及相应的和
\begin{align*}
v_{\Delta}=\bigvee_f{\left( x_0,x_1,\cdots ,x_n \right)}=\sum_{i=1}^n{|f(x_i)}-f(x_{i-1})|,
\end{align*}
称之为 $f(x)$ 在 $[a,b]$ 上的\textbf{变差}; 

作
\begin{align*}
\bigvee_a^b(f) = \sup\{v_\Delta: \Delta \text{ 为 } [a,b] \text{ 的任一分划}\},
\end{align*}
并称它为 $f(x)$ 在 $[a,b]$ 上的\textbf{全变差}.

若
\begin{align*}
\bigvee_a^b(f) < +\infty,
\end{align*}
则称 $f(x)$ 是 $[a,b]$ 上的\textbf{有界变差函数}(即全体变差形成有界数集), 其全体记为 $BV([a,b])$.
\end{definition}
\begin{remark}
由定义我们立即可得下述简单事实:
\begin{enumerate}[(1)]
\item 设 $f\in BV([a,b])$, 则 $f(x)$ 在 $[a,b]$ 上是有界函数;
\item $BV([a,b])$ 构成一个线性空间.
\end{enumerate}
\end{remark}
\begin{remark}
\textbf{有界变差函数与曲线的可求长.}

设 $f(x)$ 是定义在 $[a,b]$ 上的函数, 对 $[a,b]$ 作分划
\begin{align*}
\Delta: a = x_0 < x_1 < \cdots < x_n = b.
\end{align*}
若把依次联结点组 $\{(x_i, f(x_i))\}$ ($i=0,1,2,\cdots,n$) 的折线长记为
\begin{align*}
l_\Delta(f) = \sum_{i=1}^n \sqrt{(x_i - x_{i-1})^2 + (f(x_i) - f(x_{i-1}))^2},
\end{align*}
则定义曲线段即点集 $\{(x,y): a\leqslant x\leqslant b,\ y=f(x)\}$ 的长度为
\begin{align*}
l_f = \sup_\Delta \{l_\Delta(f)\} \quad (\text{对一切分划取上确界}).
\end{align*}
若 $l_f < +\infty$, 则称 $f(x)$ 在 $[a,b]$ 上的曲线弧段是\textbf{可求长的}.

现在从有界变差的角度来看. 因为对 $[a,b]$ 的任一分划 $\Delta: a=x_0<x_1<\cdots<x_n=b$, 我们有
\begin{align*}
|f(x_i) - f(x_{i-1})| &\leqslant \sqrt{(x_i - x_{i-1})^2 + (f(x_i) - f(x_{i-1}))^2} \\
&\leqslant (x_i - x_{i-1}) + |f(x_i) - f(x_{i-1})|,
\end{align*}
对 $i=1,2,\cdots,n$ 求和, 可得
\begin{align*}
\sum_{i=1}^n |f(x_i) - f(x_{i-1})| \leqslant l_\Delta(f) \leqslant (b-a) + \sum_{i=1}^n |f(x_i) - f(x_{i-1})|,
\end{align*}
所以 $[a,b]$ 上 $f(x)$ 之弧段可求长与 $f(x)$ 在 $[a,b]$ 上有界变差是等价的, 这可以说是有界变差函数的几何意义.
\end{remark}

\begin{proposition}\label{proposition:常见的一些有界变差函数}
设$f$是$[a,b]$上的实值函数,
\begin{enumerate}[(1)]
\item\label{proposition:常见的一些有界变差函数-1} 若$f$是单调函数,则$f\in BV[a,b]$.

\item\label{proposition:常见的一些有界变差函数-2} 若 $f(x)$ 在 $[a,b]$ 上满足Lipschitz条件
\begin{align*}
|f(x)-f(x')| \leqslant L|x-x'|, \quad \forall x,x'\in[a,b],
\end{align*}
其中 $L\geqslant 0$ 为常数, 则$f\in BV[a,b]$.

\item\label{proposition:常见的一些有界变差函数-3} 设 $f(x)$ 在 $[a,b]$ 上处处可导, 且导数一致有界: $\exists M>0$, 使得
\begin{align*}
|f'(x)| \leqslant M, \quad \forall x\in[a,b],
\end{align*}
则$f\in BV[a,b]$.
\end{enumerate}
\end{proposition}
\begin{proof}
\begin{enumerate}[(1)]
\item 不妨设 $f(x)$ 在 $[a,b]$ 上单调递增, 则对 $[a,b]$ 的任一分割 $\Delta = \{x_0, x_1, \cdots, x_n\}$, 有
\begin{align*}
\bigvee_f(x_0, \cdots, x_n) = \sum_{i=1}^n |f(x_i)-f(x_{i-1})| 
= \sum_{i=1}^n [f(x_i)-f(x_{i-1})] = f(b)-f(a).
\end{align*}
从而 $\bigvee_a^b(f) = f(b)-f(a)$. 因此, $f\in BV[a,b]$.

\item 对 $[a,b]$ 的任一分割 $\Delta = \{x_0, x_1, \cdots, x_n\}$, 有
\begin{align*}
\bigvee_f(x_0, \cdots, x_n) = \sum_{i=1}^n |f(x_i)-f(x_{i-1})| 
\leqslant \sum_{i=1}^n L|x_i-x_{i-1}| = L(b-a).
\end{align*}
从而 $\bigvee_a^b(f) = L(b-a)$. 因此, $f\in BV[a,b]$.

\item 对 $[a,b]$ 的任一分割 $\Delta = \{x_0, x_1, \cdots, x_n\}$, 有
\begin{align*}
\bigvee_f(x_0, \cdots, x_n) = \sum_{i=1}^n |f(x_i)-f(x_{i-1})|.
\end{align*}
由微分中值定理, $f(x_i)-f(x_{i-1}) = f'(\xi_i)(x_i-x_{i-1})$, $\exists \xi_i\in(x_{i-1},x_i)$. 从而
\begin{align*}
\bigvee_f(x_0, \cdots, x_n) \leqslant \sum_{i=1}^n |f'(\xi_i)||x_i-x_{i-1}| \leqslant \sum_{i=1}^n M|x_i-x_{i-1}| = M(b-a).
\end{align*}
从而 $\bigvee_a^b(f) = M(b-a)$. 因此, $f\in BV[a,b]$.
\end{enumerate}
\end{proof}

\begin{example}
设
\begin{align*}
f(x) =
\begin{cases}
x\cos\frac{\pi}{x}, & x\in(0,1], \\
0, & x=0,
\end{cases}
\end{align*}
则 $f(x)$ 显然在 $[0,1]$ 上连续, 但不是有界变差函数.
\end{example}
\begin{remark}
这个例子表明连续函数不一定是有界变差函数.
\end{remark}
\begin{proof}
对 $[0,1]$ 作如下分割
\begin{align*}
\Delta = \{x_0, x_1, \cdots, x_n\} = \left\{0, \frac{1}{n}, \frac{1}{n-1}, \cdots, \frac{1}{2}, 1\right\},
\end{align*}
则
\begin{align*}
f\left(\frac{1}{k}\right) = \frac{1}{k}\cos k\pi = (-1)^k\frac{1}{k}, \quad k=1,\cdots,n.
\end{align*}
从而
\begin{align*}
\bigvee_f(x_0, \cdots, x_n) &= \left|f\left(\frac{1}{n}\right)-f(0)\right|+\left|f\left(\frac{1}{n-1}\right)-f\left(\frac{1}{n}\right)\right|+\cdots+\left|f(1)-f\left(\frac{1}{2}\right)\right| \\
&= \frac{1}{n} + \left(\frac{1}{n-1}+\frac{1}{n}\right) + \left(\frac{1}{n-2}+\frac{1}{n-1}\right) + \cdots + \left(1+\frac{1}{2}\right) \\
&= 2\sum_{k=1}^n \frac{1}{k} - 1.
\end{align*}
令 $n\to\infty$, 得 $\bigvee_a^b(f) = +\infty$. 故 $f\notin BV[a,b]$.
\end{proof}

\begin{proposition}[有界变差函数的性质]\label{proposition:有界变差函数的性质}
设 $f,g\in BV[a,b]$, 则
\begin{enumerate}[(1)]
\item\label{proposition:有界变差函数的性质-1} $f$ 是有界函数;

\item\label{proposition:有界变差函数的性质-2} $kf\in BV[a,b]$, 且 $\bigvee_a^b(kf)\leqslant|k|\bigvee_a^b(f)$;

\item\label{proposition:有界变差函数的性质-3} $f+g\in BV[a,b]$, 且 $\bigvee_a^b(f+g)\leqslant\bigvee_a^b(f)+\bigvee_a^b(g)$;

\item\label{proposition:有界变差函数的性质-4} $fg\in BV[a,b]$;

\item\label{proposition:有界变差函数的性质-5} 对 $\forall c\in[a,b]$ 都有 $\bigvee_a^b(f)=\bigvee_a^c(f)+\bigvee_c^b(f)$.
\end{enumerate}
\end{proposition}
\begin{proof}
\begin{enumerate}[(1)]
\item 假设 $f$ 无界, 则存在 $\{x_n\}\subset[a,b]$ 使得 $|f(x_n)|\to+\infty$. 由于 $f\in BV[a,b]$, 故 $\bigvee_a^b(f)<+\infty$. 对 $\forall n\in\mathbb{N}$, 作分割 $\Delta_n=\{a,x_n,b\}$. 则
\begin{align*}
\bigvee_f(a,x_n,b) = |f(x_n)-f(a)|+|f(b)-f(x_n)| 
\geqslant 2|f(x_n)|-|f(a)|-|f(b)|.
\end{align*}
从而, 对任意的 $n\in\mathbb{N}$ 有
\begin{align*}
|f(x_n)| \leqslant \frac{1}{2}\left[\bigvee_f(a,x_n,b)+|f(a)|+|f(b)|\right] \leqslant \frac{1}{2}\left[\bigvee_a^b(f)+|f(a)|+|f(b)|\right]<+\infty,
\end{align*}
矛盾.

\item 容易验证.

\item 容易验证.

\item  $f,g\in BV[a,b]$, 由 (1) 知, $f,g$ 有界, 故存在 $M_1,M_2>0$, 使得
\begin{align*}
|f(x)|\leqslant M_1,\quad |g(x)|\leqslant M_2,\quad \forall x\in[a,b].
\end{align*}
对 $[a,b]$ 的任一分割 $\Delta=\{x_0,\cdots,x_n\}$, 有
\begin{align*}
\bigvee_{fg}(x_0,\cdots,x_n) &= \sum_{i=1}^n |(fg)(x_i)-(fg)(x_{i-1})| 
= \sum_{i=1}^n |f(x_i)g(x_i)-f(x_{i-1})g(x_{i-1})| \\
&\leqslant \sum_{i=1}^n \left[|f(x_i)g(x_i)-f(x_i)g(x_{i-1})|+|f(x_i)g(x_{i-1})-f(x_{i-1})g(x_{i-1})|\right] \\
&\leqslant M_1\sum_{i=1}^n |g(x_i)-g(x_{i-1})|+M_2\sum_{i=1}^n |f(x_i)-f(x_{i-1})| \\
&\leqslant M_1\bigvee_a^b(g)+M_2\bigvee_a^b(f).
\end{align*}
从而
\begin{align*}
\bigvee_a^b(fg)\leqslant M_1\bigvee_a^b(g)+M_2\bigvee_a^b(f).
\end{align*}

\item 对 $[a,c]$ 的任一分割 $\{x_0,\cdots,x_n\}$, 以及 $[c,b]$ 的任一分割 $\{x_0',\cdots,x_m'\}$, 都可合并为 $[a,b]$ 的一个分割
\begin{align*}
a=x_0<x_1<\cdots<x_n=c=x_0'<x_1'<\cdots<x_m'=b.
\end{align*}
则
\begin{align*}
\bigvee_f(x_0,\cdots,x_n)+\bigvee_f(x_0',\cdots,x_m') &= \sum_{i=1}^n |f(x_i)-f(x_{i-1})|+\sum_{j=1}^m |f(x_j')-f(x_{j-1}')| \\
&= \bigvee_f(x_0,\cdots,x_n,x_1',\cdots,x_m')
\leqslant \bigvee_a^b(f).
\end{align*}
对左端所有分割取上确界, 得 $\bigvee_a^c(f)+\bigvee_c^b(f)\leqslant\bigvee_a^b(f)$.

另一方面, 对 $\forall\varepsilon>0$, 存在 $[a,b]$ 的一个分割 $\{x_0,\cdots,x_n\}$ 使得
\begin{align*}
\bigvee_f(x_0,\cdots,x_n)>\bigvee_a^b(f)-\varepsilon.
\end{align*}
设 $c\in(x_{k-1},x_k]$, 则 $\{x_0,\cdots,x_{k-1},c\}$ 和 $\{c,x_k,\cdots,x_n\}$ 分别是 $[a,c]$ 和 $[c,b]$ 的分割. 注意到, 对分割增加分点后, 新的变差不会减小, 故
\begin{align*}
\bigvee_f(x_0,\cdots,x_n) &\leqslant \bigvee_f(x_0,\cdots,x_{k-1},c,x_k,\cdots,x_n) \\
&= \bigvee_f(x_0,\cdots,x_{k-1},c)+\bigvee_f(c,x_k,\cdots,x_n) \\
&\leqslant \bigvee_a^c(f)+\bigvee_c^b(f).
\end{align*}
从而 $\bigvee_a^b(f)-\varepsilon\leqslant\bigvee_a^c(f)+\bigvee_c^b(f)$, 再由 $\varepsilon$ 的任意性, $\bigvee_a^b(f)\leqslant\bigvee_a^c(f)+\bigvee_c^b(f)$.
\end{enumerate}
\end{proof}

\begin{definition}[变差函数]
设 $f(x)$ 为 $[a,b]$ 上的有界变差函数, 则对每一个 $x\in[a,b]$, 由\hyperref[proposition:有界变差函数的性质]{有界变差函数的性质\ref{proposition:有界变差函数的性质-5}}知, $f$ 在 $[a,x]$ 上也是有界变差的, 从而 $\bigvee_a^x(f)$ 是 $[a,b]$ 上的实值函数, 称为 $f$ 的\textbf{变差函数}.
\end{definition}
\begin{remark}
显然, $\bigvee_a^x(f)$ 是单调递增函数.
\end{remark}

\begin{theorem}[Jordan(约当)分解定理]\label{theorem:实变函数--Jordan(约当)分解定理}
$f(x)$ 是 $[a,b]$ 上的有界变差函数, 当且仅当
\begin{align*}
f(x) = g(x) - h(x),
\end{align*}
其中, $g(x), h(x)$ 为 $[a,b]$ 上的单调递增函数.
\end{theorem}
\begin{remark}
我们称$f(x)=g(x)-h(x)$为\textbf{标准}$\mathbf{Jordan}$分解,其中
\begin{align*}
g(x) = \frac{1}{2}\left(\bigvee_a^x(f) + f(x)\right), \quad h(x) = \frac{1}{2}\left(\bigvee_a^x(f) - f(x)\right).
\end{align*}
\end{remark}
\begin{proof}

$\Longleftarrow$:由\rrefpro{proposition:常见的一些有界变差函数}{proposition:常见的一些有界变差函数-1}和\hyperref[proposition:有界变差函数的性质]{有界变差函数的性质\ref{proposition:有界变差函数的性质-3}}立得.

$\Longrightarrow$:设$f\in BV[a,b]$,令
\begin{align*}
g(x) = \frac{1}{2}\left(\bigvee_a^x(f) + f(x)\right), \quad h(x) = \frac{1}{2}\left(\bigvee_a^x(f) - f(x)\right),
\end{align*}
则 $f(x) = g(x) - h(x)$. 下面证 $g$ 单调递增. 设 $x_1 < x_2$, 由\hyperref[proposition:有界变差函数的性质]{有界变差函数的性质\ref{proposition:有界变差函数的性质-5}}得
\begin{align*}
f(x_1) - f(x_2) \leqslant \vee_f(x_1,x_2) \leqslant \bigvee_{x_1}^{x_2}(f) = \bigvee_a^{x_2}(f) - \bigvee_a^{x_1}(f).
\end{align*}
从而, $\bigvee_a^{x_1}(f) + f(x_1) \leqslant \bigvee_a^{x_2}(f) + f(x_2)$, 于是 $g(x_1) \leqslant g(x_2)$. 类似可证, $h$ 也单调递增.
\end{proof}

\begin{corollary}
设 $f(x)$ 为 $[a,b]$ 上的有界变差函数,
\begin{enumerate}[(1)]
\item $f(x)$ 不连续点的全体至多是可列集;
\item $f(x)$ 在 $[a,b]$ 上 R-可积, 从而 L-可积;
\item $f(x)$ 在 $[a,b]$ 上几乎处处可导, 且 $f'(x)$ 是 L-可积的.
\end{enumerate}
\end{corollary}
\begin{proof}

\end{proof}

\begin{theorem}
设 $f(x)$ 是 $[a,b]$ 上的有界变差函数, 则
\begin{align*}
\bigvee_a^x(f) \text{ 在 } [a,b] \text{ 上右连续 (左连续)} \iff f(x) \text{ 在 } [a,b] \text{ 上右连续 (左连续)}.
\end{align*}
\end{theorem}
\begin{proof}

$\Longrightarrow$: 设 $\bigvee_a^x(f)$ 在 $[a,b]$ 上右连续, 任取 $x_0\in[a,b)$, 对 $\forall x\in(x_0,b]$, 有
\begin{align*}
|f(x)-f(x_0)| = \vee_f(x_0,x) \leqslant \bigvee_{x_0}^x(f) = \bigvee_a^x(f) - \bigvee_a^{x_0}(f).
\end{align*}
从而, $\bigvee_a^x(f)$ 右连续 $\Rightarrow f(x)$ 右连续.

$\Longleftarrow$: 设 $f(x)$ 在 $[a,b]$ 上右连续, 任取 $x_0\in[a,b)$, 则对 $\forall\varepsilon>0$, $\exists\delta>0$ ($\delta\leqslant b-x_0$), 当 $x\in(x_0,x_0+\delta)$ 时, 有
\begin{align*}
|f(x)-f(x_0)| < \frac{\varepsilon}{2}.
\end{align*}
取 $[x_0,x_0+\delta]$ 的一个分割 $\Delta = \{t_0,\cdots,t_n\}$, 其中
\begin{align*}
x_0 = t_0 < t_1 < \cdots < t_n = x_0+\delta
\end{align*}
使得
\begin{align}
\sum_{i=1}^n |f(t_i)-f(t_{i-1})| > \bigvee_{x_0}^{x_0+\delta}(f) - \frac{\varepsilon}{2}. \label{eq:::-hr82h39rhswufh8jiojshiojfw812}
\end{align}
又 $\Delta_1 = \{t_1,\cdots,t_n\}$ 是 $[t_1,x_0+\delta]$ 的一个分割, 故
\begin{align}
\sum_{i=2}^n |f(t_i)-f(t_{i-1})| \leqslant \bigvee_{t_1}^{x_0+\delta}(f). \label{eq:::-hr82h39rhswufh8jiojshiojfw813}
\end{align}
由式 \eqref{eq:::-hr82h39rhswufh8jiojshiojfw812} 和 \eqref{eq:::-hr82h39rhswufh8jiojshiojfw813} 得
\begin{align*}
\bigvee_{x_0}^{t_1}(f) &= \bigvee_{x_0}^{x_0+\delta}(f) - \bigvee_{t_1}^{x_0+\delta}(f) 
\leqslant \sum_{i=1}^n |f(t_i)-f(t_{i-1})| + \frac{\varepsilon}{2} - \sum_{i=2}^n |f(t_i)-f(t_{i-1})| \\
&= |f(t_1)-f(t_0)| + \frac{\varepsilon}{2} < \varepsilon.
\end{align*}
从而, 当 $x\in[x_0,t_1]$ 时, 有
\begin{align*}
\bigvee_a^x(f) - \bigvee_a^{x_0}(f) = \bigvee_{x_0}^x(f) \leqslant \bigvee_{x_0}^{t_1}(f) < \varepsilon.
\end{align*}
因此, $\bigvee_a^x(f)$ 在点 $x_0$ 右连续.
\end{proof}

\begin{proposition}\label{proposition:有界变差函数必几乎处处可微}
若 $f\in BV([a,b])$, 则 $f(x)$ 几乎处处可微, 且
\begin{align*}
\frac{\mathrm{d}}{\mathrm{d}x}\left(\bigvee_a^x(f)\right) = |f'(x)|, \quad \text{a.e. } x\in[a,b].
\end{align*}
\end{proposition}
\begin{proof}
(i) 根据全变差的定义可知, 对任给的 $\varepsilon>0$, 存在分划
\begin{align*}
\Delta: a = x_0 < x_1 < \cdots < x_k = b,
\end{align*}
使得
\begin{align}
\bigvee_a^b(f) - \sum_{i=1}^k |f(x_i)-f(x_{i-1})| < \varepsilon. \label{eq:5.10}
\end{align}
我们可以做出 $[a,b]$ 上的函数 $g(x)$, 使得
\begin{align*}
g(x) =
\begin{cases}
f(x)+c_i, & f(x_i)\geqslant f(x_{i-1}), \\
-f(x)+c_i', & f(x_i)<f(x_{i-1})
\end{cases}
\quad \left(
\begin{aligned}
x&\in[x_{i-1},x_i], \\
i&=1,2,\cdots,k
\end{aligned}
\right),
\end{align*}
其中 $c_i,c_i'$ ($i=1,2,\cdots,k$) 是常数.

实际上, 首先, 对 $x\in[a,x_1]$, 令
\begin{align*}
g(x) =
\begin{cases}
f(x)-f(a_0), & f(x_1)\geqslant f(a_0), \\
-f(x)+f(a_0), & f(x_1)<f(a_0).
\end{cases}
\end{align*}
其次, 用归纳法, 若在 $[a,x_i]$ ($i<k$) 上已定义了 $g(x)$, 则对 $x\in(x_i,x_{i+1}]$, 定义
\begin{align*}
g(x) =
\begin{cases}
f(x)+(g(x_i)-f(x_i)), & f(x_{i+1})\geqslant f(x_i), \\
-f(x)+(g(x_i)+f(x_i)), & f(x_{i+1})<f(x_i).
\end{cases}
\end{align*}
对如此做成的 $g(x)$, 易知对每个 $[x_{i-1},x_i]$, $g(x)-f(x)$ 或 $g(x)+f(x)$ 是常数, 且 $|g'(x)|=|f'(x)|$, a.e. $x\in[a,b]$. 又由 \eqref{eq:5.10} 式可得
\begin{align*}
\bigvee_a^b(f) - g(b) < \varepsilon,
\end{align*}
以及 $\bigvee_a^x(f)-g(x)$ 是 $[a,b]$ 上的递增函数.

(ii) 由 (i) 知, 对 $\varepsilon=\frac{1}{2}^n$, 存在 $[a,b]$ 上的函数列 $\{g_n(x)\}$, 使得
\begin{align*}
\bigvee_a^b(f) - g_n(b) < \frac{1}{2^n}, \quad |g_n'(x)| = |f'(x)|, \quad \text{a.e. } x\in[a,b].
\end{align*}
这说明
\begin{align*}
\sum_{n=1}^\infty \left(\bigvee_a^x(f) - g_n(x)\right) < +\infty, \quad x\in[a,b].
\end{align*}
引用\hyperref[theorem:实变函数--Fubini 逐项微分定理]{Fubini 逐项微分定理}, 可知
\begin{align*}
\sum_{n=1}^\infty \left(\frac{\mathrm{d}}{\mathrm{d}x}\bigvee_a^x(f) - g_n'(x)\right) < +\infty, \quad \text{a.e. } x\in[a,b].
\end{align*}
由于 $|g_n'(x)|=|f'(x)|$, a.e. $x\in[a,b]$, 且 $\frac{\mathrm{d}}{\mathrm{d}x}\bigvee_a^x(f)\geqslant 0$, a.e. $x\in[a,b]$, 故我们有
\begin{align*}
\frac{\mathrm{d}}{\mathrm{d}x}\bigvee_a^x(f) = |f'(x)|, \quad \text{a.e. } x\in[a,b].
\end{align*}
\end{proof}

\begin{proposition}\label{proposition:实变函数--弧长公式1}
若 $f\in BV([a,b])$, 则弧长
\begin{align*}
l_f \geqslant \int_a^b \sqrt{1+(f'(x))^2}\mathrm{d}x.
\end{align*}
\end{proposition}
\begin{proof}
实际上, 记 $l_f(x)$ 是在 $[a,x]$ 上曲线 $y=f(x)$ 的弧长 ($a\leqslant x\leqslant b$), 易知对 $a\leqslant x<y\leqslant b$, 有
\begin{align*}
l_f(y)-l_f(x) \geqslant \sqrt{(y-x)^2+(f(y)-f(x))^2} > 0,
\end{align*}
即 $l_f(x)$ 在 $[a,b]$ 上递增. 由于
\begin{align*}
\frac{l_f(y)-l_f(x)}{y-x} \geqslant \sqrt{1+\left(\frac{f(y)-f(x)}{y-x}\right)^2},
\end{align*}
故得
\begin{align*}
l_f'(x) \geqslant \sqrt{1+(f'(x))^2}, \quad \text{a.e. } x\in[a,b].
\end{align*}
从而我们有
\begin{align*}
l_f(b) \geqslant \int_a^b l_f'(x)\mathrm{d}x \geqslant \int_a^b \sqrt{1+(f'(x))^2}\mathrm{d}x.
\end{align*}
\end{proof}

\begin{proposition}[Poisson公式]
设 $f\in C(\mathbb{R})\cap L(\mathbb{R})$. 若 $f\in BV(\mathbb{R})$, 则
\begin{align*}
\sum_{n=-\infty}^{+\infty} f(2n\pi) = \frac{1}{2\pi}\sum_{n=-\infty}^{+\infty} \hat{f}(n),
\end{align*}
其中 $\hat{f}(n)$ 表示 Fourier 系数.
\end{proposition}
\begin{proof}
作函数 $f^*(x)=\sum_{n=-\infty}^{+\infty} f(x+2n\pi)$ ($x\in\mathbb{R}$), 以及 $V(x)=\bigvee_{-\infty}^x(f)$. 易知 $f^*(x)$ 在 $\mathbb{R}$ 上几乎处处有定义且可积, 又存在极限 $V(x)\to V$ ($x\to+\infty$). 我们有
\begin{align}\label{eq:::290jr98wjef}
\int_0^{2\pi} f^*(x)\mathrm{d}x = \int_{-\infty}^{+\infty} f(x)\mathrm{d}x.
\end{align}
取 $x_0\in\mathbb{R}$, 使得 $f^*(x_0)$ 有定义, 则对 $x_0\leqslant x\leqslant x_0+2\pi$, 可知
\begin{align*}
|f(x+2n\pi)-f(x_0+2n\pi)| \leqslant \bigvee_{x_0}^{x+2n\pi}(f) - \bigvee_{-\infty}^{x_0+2n\pi}(f) 
\leqslant \bigvee_{-\infty}^{x_0+2(n+1)\pi}(f) - \bigvee_{-\infty}^{x_0+2n\pi}(f), 
\end{align*}
\begin{align*}
\bigvee_{-\infty}^{+\infty}\left(\bigvee_{-\infty}^{x_0+2(n+1)\pi}(f) - \bigvee_{-\infty}^{x_0+2n\pi}(f)\right) = V < +\infty.
\end{align*}
由此得 $f^*\in C([x_0,x_0+2\pi])$, 从而有 $f^*\in C(\mathbb{R})$.

对于 $0=x_0<x_1<\cdots<x_k=2\pi$, 有
\begin{align*}
\sum_{i=0}^{k-1} |f^*(x_{i+1})-f^*(x_i)|
&= \sum_{i=0}^{k-1} \left|\sum_{n=-\infty}^{+\infty} f(x_{i+1}+2n\pi) - \sum_{n=-\infty}^{+\infty} f(x_i+2n\pi)\right| \\
&\leqslant \sum_{i=0}^{k-1} \sum_{n=-\infty}^{+\infty} |f(x_{i+1}+2n\pi)-f(x_i+2n\pi)| \leqslant V.
\end{align*}
应用公式\eqref{eq:::290jr98wjef}于 $\mathrm{e}^{-\mathrm{i}nx}f(x)$, 则得
\begin{align*}
\frac{1}{2\pi}\int_0^{2\pi} \mathrm{e}^{-\mathrm{i}nx}f^*(x)\mathrm{d}x = \frac{1}{2\pi}\int_{-\infty}^{+\infty} \mathrm{e}^{-\mathrm{i}nx}f(x)\mathrm{d}x = \frac{1}{2\pi}\hat{f}(n).
\end{align*}
最后 (根据 Jordan-Dirichlet 定理), 我们有
\begin{align*}
\sum_{n=-\infty}^{\infty} f(2n\pi) = f^*(0) = \sum_{n=-\infty}^{\infty} \frac{1}{2\pi}\int_0^{2\pi} \mathrm{e}^{-\mathrm{i}nx}f^*(x)\mathrm{d}x = \frac{1}{2\pi}\sum_{n=-\infty}^{\infty} \hat{f}(n).
\end{align*}
\end{proof}




\end{document}